\uuid{db2e}
\titre{Matrice d'une application linéaire/Noyau/Image/Espaces supplémentaires}
\auteur{hanani ; mzari}
\organisation{exo7}
\licence{CC BY-NC-SA 4.0}
\datecreate{2020-01-20}
\niveau{L1}
\module{Algèbre}
\chapitre{Matrice}
\sousChapitre{Matrice et application linéaire}
\difficulte{2}
\theme{Matrice d'une application linéaire, Noyau, Image, Espaces supplémentaires}
\exo7qcmid{lille-2.matrice-application.012}

\contenu{
\begin{question}
\qclassification{108.03}
\qtags{motcle=Matrice d'une application linéaire/Noyau/Image/Espaces supplémentaires}
On considère $\Rr^3$ muni de sa base canonique notée ${\cal {B}}$ et
$f$ l'application linéaire définie  par :
$$\begin{array}{rccc}f:&\Rr^3 &\to& \Rr^3\\
& (x,y,z)&\to &(x+z,2x+2z,-x-z).  \end{array}$$
Soit $ u_1=(0,1,0), u_2=(1,2,-1)$ et $ u_3=(1,0,0)$.
Quelles sont les assertions vraies ?
\begin{answers}
\good{$\{u_1,u_2,u_3\}$ est une base de $\Rr^3$}.
\good{$\{u_1 , u_2\}$  est une base de $\ker f$.}
\bad{$\ker f$ et $\Im f$ sont supplémentaires dans $\Rr^3$}.
\good{La matrice de $f$ dans la base  $\{u_1,u_2,u_3\}$ est :
$\left(\begin{array}{rcc}
0&0&0\\
0&0&1\\
0&0&0\\
\end{array}\right)$}.
\end{answers}
\begin{explanations} $\{u_1,u_2,u_3\}$ est une base de $\Rr^3$,
$\{u_1 , u_2\}$  est une base de $\ker f$ et  $\{u_2\}$  est une base de $\Im f$. Comme $ \ker f \cap \Im f$ est non nul, $\ker f $ et $\Im f$ ne sont pas supplémentaires dans $\Rr^3$. Comme $f(u_1)=f(u_2)=0$ et $f(u_3)=u_2$, la matrice dans la base
$\{u_1,u_2,u_3\}$ est :
$$\left(\begin{array}{rcc}
0&0&0\\
0&0&1\\
0&0&0\end{array}\right).$$
\end{explanations}
\end{question}
}
